% =============================================================================
% SECTION: Schede Personaggio — tutti e 20 i PG
% @rpg-schema: fitd:Character a owl:Class ;
%   rdfs:label "Personaggio" .
% @rpg-schema: fitd:Bond a owl:ObjectProperty ;
%   rdfs:label "Legame" ;
%   rdfs:domain fitd:Character ;
%   rdfs:range fitd:Character .
% =============================================================================

\chapter{Schede Personaggio — Casa Doria}
\label{ch:pg-doria}

% ---------------------------------------------------------------------------
% PG 1: Ottaviano Doria
% @rpg-schema: pg:OttavianoDoria a fitd:Character ;
%   rdfs:label "Ottaviano Doria" ;
%   fitd:crew genova1507:CasaDoria ;
%   rpg:role "Capofamiglia / Ammiraglio" ;
%   fitd:insight 7 ; fitd:prowess 6 ; fitd:resolve 6 .
% ---------------------------------------------------------------------------

\charheader{doria}{Ottaviano Doria}{Il Capofamiglia / L'Ammiraglio}%
  {Il mare non chiede permesso. Nemmeno io.}

\begin{multicols}{2}

\subsection*{Background \& Motivazione}
% @rpg-schema: pg:OttavianoDoria rpg:background
%   "Cinquantadue anni, cicatrici da battaglie navali, comandante dai vent'anni." .

Cinquantadue anni, cicatrici da combattimenti nelle acque greche quando aveva vent'anni.
È il volto ufficiale della famiglia a Genova — quello che i francesi invitano
ai banchetti e che sorride mentre pianifica la loro rovina.

\textbf{\textcolor{DoriaBlue}{Motivazione:}} \textit{Vuole rivedere Genova libera prima di
morire — non per ideologia, ma per orgoglio. I Doria governano il mare
da trecento anni. Non tollera che un re straniero dica loro dove attraccare.}

\columnbreak

\subsection*{Stress \& Trauma}
\clocktrack[DoriaBlue]{STRESS}{0}{9}
\textit{\small Traumi: □ Freddo \ □ Duro \ □ Ossessionato \ □ Paranoico \ □ Spezzato}

\end{multicols}

\attrblock{doria}{INSIGHT — Mente}
% @rpg-schema: pg:OttavianoDoria fitd:hasAction
%   [ fitd:actionName "Caccia" ; fitd:rating 2 ] ,
%   [ fitd:actionName "Studio" ; fitd:rating 1 ] ,
%   [ fitd:actionName "Osserva" ; fitd:rating 2 ] ,
%   [ fitd:actionName "Ingegno" ; fitd:rating 0 ] .
\statrow{Caccia (Hunt)}   {2}{4}{Individua bersagli, traccia rotte, legge le mappe militari}
\statrow{Studio (Study)}  {1}{4}{Analisi tattica, storia navale, intenzioni nemiche}
\statrow{Osserva (Survey)}{2}{4}{Valuta situazioni, scansiona l'orizzonte politico}
\statrow{Ingegno (Tinker)}{0}{4}{Meccanismi, artiglieria, strutture navali}

\attrblock{doria}{PROWESS — Corpo}
% @rpg-schema: pg:OttavianoDoria fitd:hasAction
%   [ fitd:actionName "Finezza" ; fitd:rating 1 ] ,
%   [ fitd:actionName "Ombra" ; fitd:rating 0 ] ,
%   [ fitd:actionName "Combattimento" ; fitd:rating 3 ] ,
%   [ fitd:actionName "Forza Bruta" ; fitd:rating 2 ] .
\statrow{Finezza (Finesse)}{1}{4}{Precisione, mani abili, duello con stile}
\statrow{Ombra (Prowl)}{0}{4}{Muoversi non visti, infiltrazione}
\statrow{Combattimento (Skirmish)}{3}{4}{Scontro diretto, spada, corpo a corpo}
\statrow{Forza Bruta (Wreck)}{2}{4}{Sfondare porte, affondare navi, distruggere}

\attrblock{doria}{RESOLVE — Spirito}
% @rpg-schema: pg:OttavianoDoria fitd:hasAction
%   [ fitd:actionName "Percezione" ; fitd:rating 0 ] ,
%   [ fitd:actionName "Comando" ; fitd:rating 3 ] ,
%   [ fitd:actionName "Relazioni" ; fitd:rating 2 ] ,
%   [ fitd:actionName "Persuasione" ; fitd:rating 1 ] .
\statrow{Percezione (Attune)}{0}{4}{Leggere atmosfere, sentire il non detto}
\statrow{Comando (Command)}{3}{4}{Guidare uomini, imporre la propria volontà}
\statrow{Relazioni (Consort)}{2}{4}{Frequentare i potenti, muoversi nei palazzi}
\statrow{Persuasione (Sway)}{1}{4}{Convincere, manipolare, negoziare}

\vspace{4pt}
% @rpg-schema: pg:OttavianoDoria fitd:hasSpecialAbility genova1507:AutoritaAmmiraglio .
% @rpg-schema: genova1507:AutoritaAmmiraglio a fitd:SpecialAbility ;
%   rdfs:label "Autorità dell'Ammiraglio" ;
%   fitd:frequency "una volta per colpo" .
\abilitybox{doria}{Autorità dell'Ammiraglio}{%
  Quando Ottaviano impartisce un ordine in una situazione di crisi, tutti i PG del
  tavolo che lo eseguono tirano con Posizione \textbf{Controllata} invece di Rischiosa.%
}{Una volta per colpo. Deve essere un ordine esplicito pronunciato in scena.}

\subsection*{Legami}
% @rpg-schema: pg:OttavianoDoria fitd:hasBond
%   [ fitd:bondTarget pg:FulvioDoria ; fitd:bondText "Luogotenente da trent'anni. Mente quando evita il mio sguardo — e lo fa troppo spesso." ] ,
%   [ fitd:bondTarget pg:BiancaDoria ; fitd:bondText "Nipote. Più intelligente di chiunque in questa stanza." ] ,
%   [ fitd:bondTarget pg:FraMatteo ; fitd:bondText "Cappellano. Lo salvai dai turchi nel '78. Mi deve qualcosa che non potrà mai restituire." ] ,
%   [ fitd:bondTarget pg:GiacomoMoro ; fitd:bondText "Contrabbandiere. Non mi fido di lui, ma è il migliore nel suo mestiere." ] .
\bondentry{Ser Fulvio}{Il mio luogotenente da trent'anni. So che mente quando evita il mio sguardo — e lo fa troppo spesso ultimamente.}
\bondentry{Bianca}{Mia nipote. Più intelligente di chiunque in questa stanza. La proteggerei con la vita, ma lei non ne ha bisogno.}
\bondentry{Fra' Matteo}{Lo salvai dai turchi nel '78. Mi deve qualcosa che non potrà mai restituire, e questo mi pesa.}
\bondentry{Giacomo il Moro}{Non mi fido di lui, ma è il migliore nel suo mestiere.}

\subsection*{Equipaggiamento}
\begin{itemize}
  \item Spada da lato di fattura spagnola — un dono di cui non ricorda l'origine
  \item Carte nautiche segrete del porto con i fondali reali (non quelli per i francesi)
  \item Sigillo ufficiale della famiglia — vale più di cento spade
  \item Fiaschetta d'argento con grappa delle Cinque Terre. Non offre mai da bere.
\end{itemize}

\begin{secretbox}
% @rpg-schema: pg:OttavianoDoria rpg:hasSecret
%   "Ha ipotecato segretamente due galee con un banchiere filofrancese." .
% @rpg-schema: pg:OttavianoDoria rpg:hasHiddenAgenda
%   "Completare il Colpo Riservato a qualsiasi costo personale. Mai le navi." .
Tre anni fa, durante una crisi di liquidità, Ottaviano ha ipotecato in segreto
due delle quattro galee familiari con un banchiere genovese di simpatie francesi.
Se la famiglia fallisce nel colpo riservato, il banchiere può confiscare le navi.
Nessun altro Doria lo sa.

\medskip
\textbf{\textcolor{SecretRed}{Agenda:}} Completare il Colpo Riservato a qualsiasi costo personale.
Se necessario, sacrificare la propria reputazione o un compagno — ma non le navi. Mai le navi.
\end{secretbox}

\newpage

% ---------------------------------------------------------------------------
% PG 2: Ser Fulvio Doria
% @rpg-schema: pg:FulvioDoria a fitd:Character ;
%   rdfs:label "Ser Fulvio Doria" ;
%   fitd:crew genova1507:CasaDoria ;
%   rpg:role "Luogotenente / Braccio Destro" .
% ---------------------------------------------------------------------------

\charheader{doria}{Ser Fulvio Doria}{Il Luogotenente / Il Braccio Destro}%
  {Un uomo che non mente non sopravvive in questa città. Io sopravvivo.}

\begin{multicols}{2}
\subsection*{Background \& Motivazione}
Quarantacinque anni, cugino di secondo grado del capofamiglia.
Ha gestito le operazioni sporche della famiglia per vent'anni.
Sa dove sono sepolti i segreti di Genova. Letteralmente.

\textbf{\textcolor{DoriaBlue}{Motivazione:}} \textit{Fedele ai Doria, con riserva mentale:
la fedeltà ha un prezzo, e qualcuno gli deve ancora molte cose.}
\columnbreak
\subsection*{Stress \& Trauma}
\clocktrack[DoriaBlue]{STRESS}{0}{9}
\textit{\small Traumi: □ Freddo \ □ Duro \ □ Ossessionato \ □ Paranoico \ □ Spezzato}
\end{multicols}

\attrblock{doria}{INSIGHT}
\statrow{Caccia}{2}{4}{Trova persone, localizza bersagli mobili, segue tracce}
\statrow{Studio}{1}{4}{Legge documenti, valuta testimonianze}
\statrow{Osserva}{3}{4}{Non sfugge nulla al suo sguardo — soprattutto le menzogne}
\statrow{Ingegno}{1}{4}{Scassinatore occasionale, meccanismi semplici}

\attrblock{doria}{PROWESS}
\statrow{Finezza}{2}{4}{Borseggio, veleno, mani veloci}
\statrow{Ombra}{3}{4}{Sorveglianza, pedinamento, sparire in una folla}
\statrow{Combattimento}{2}{4}{Coltello, scontri ravvicinati sporchi}
\statrow{Forza Bruta}{0}{4}{Non è la sua specialità}

\attrblock{doria}{RESOLVE}
\statrow{Percezione}{1}{4}{Sente le tensioni in una stanza}
\statrow{Comando}{1}{4}{Impartisce ordini ai subalterni}
\statrow{Relazioni}{2}{4}{Frequenta taverne, bordelli, e i posti dove si parla}
\statrow{Persuasione}{3}{4}{Convince, minaccia, ricatta con eleganza}

\abilitybox{doria}{L'Informatore nell'Ombra}{%
  Una volta per sessione, Fulvio può dichiarare di avere già predisposto
  una risorsa utile nella scena attuale — un testimone corrotto, una via di
  fuga, un documento compromettente — senza che fosse annunciato prima.%
}{Nessun tiro. La risorsa esiste. Il GM determina come entra in gioco.}

\subsection*{Legami}
% @rpg-schema: pg:FulvioDoria fitd:hasBond
%   [ fitd:bondTarget pg:OttavianoDoria ; fitd:bondText "L'Ammiraglio mi crede fedele oltre ogni dubbio. Ha ragione — quasi." ] ,
%   [ fitd:bondTarget pg:GiacomoMoro ; fitd:bondText "Siamo cugini veri, di sangue. Lo proteggerei a dispetto di qualsiasi ordine." ] .
\bondentry{Ottaviano}{L'Ammiraglio mi crede fedele oltre ogni dubbio. Ha ragione — quasi. C'è una cosa che non sa.}
\bondentry{Bianca}{La nipote mi spaventa. Vede troppo. Mi ha guardato una volta in un modo che mi ha fatto capire che sa qualcosa.}
\bondentry{Fra' Matteo}{Lo uso per i messaggi più delicati. Sa che se parlasse, non rivedrebbe il suo convento.}
\bondentry{Giacomo}{Siamo cresciuti insieme nei caruggi. È mio cugino vero, di sangue. Lo proteggerei a dispetto di qualsiasi ordine.}

\subsection*{Equipaggiamento}
\begin{itemize}
  \item Stiletto nascosto nello stivale destro
  \item Tre documenti falsi con identità diverse, pronti all'uso
  \item Lista criptata di debitori e creditori — solo lui sa leggerla
  \item Chiave di una casa nel sestiere di Soziglia che ufficialmente non esiste
\end{itemize}

\begin{secretbox}
% @rpg-schema: pg:FulvioDoria rpg:hasSecret
%   "È in contatto con gli Spinola: crede di usarli, ma l'ultima informazione passata era vera e ha danneggiato i Doria." .
Fulvio è in contatto con gli Spinola. Li ha convinti di essere un informatore
quando in realtà filtra ciò che vuole. Ma l'ultima informazione che ha passato
era vera — e ha danneggiato i Doria.

\medskip
\textbf{\textcolor{SecretRed}{Agenda:}} Garantirsi una via di uscita personale qualunque sia
l'esito della rivolta. Ha già trattato la propria sopravvivenza con una terza parte.
\end{secretbox}

\newpage

% ---------------------------------------------------------------------------
% PG 3: Bianca Doria
% @rpg-schema: pg:BiancaDoria a fitd:Character ;
%   rdfs:label "Bianca Doria" ;
%   fitd:crew genova1507:CasaDoria ;
%   rpg:role "La Stratega / L'Erede Ombra" .
% ---------------------------------------------------------------------------

\charheader{doria}{Bianca Doria}{La Stratega / L'Erede Ombra}%
  {Mio nonno comanda le galee. Io comando mio nonno.}

\begin{multicols}{2}
\subsection*{Background \& Motivazione}
Ventisette anni, nipote di Ottaviano. Non avrebbe dovuto contare nulla —
è una donna, in un mondo di uomini. Invece è lei che ha pianificato
il Colpo Riservato e identificato la rotta del convoglio francese.
Ottaviano lo sa, e la chiama «la mia nipotina» davanti agli altri.

\textbf{\textcolor{DoriaBlue}{Motivazione:}} \textit{Dimostrare che il casato Doria può
prosperare quando comandano le persone giuste, non solo quelle con la spada più lunga.}
\columnbreak
\subsection*{Stress \& Trauma}
\clocktrack[DoriaBlue]{STRESS}{0}{9}
\textit{\small Traumi: □ Freddo \ □ Duro \ □ Ossessionato \ □ Paranoico \ □ Spezzato}
\end{multicols}

\attrblock{doria}{INSIGHT}
\statrow{Caccia}{1}{4}{Individua pattern, cerca informazioni metodicamente}
\statrow{Studio}{4}{4}{Analisi, lingue (latino, francese, arabo), cartografia}
\statrow{Osserva}{3}{4}{Legge stanze, persone, situazioni politiche}
\statrow{Ingegno}{2}{4}{Codici, meccanismi, trappole logistiche}

\attrblock{doria}{PROWESS}
\statrow{Finezza}{2}{4}{Scrive, disegna, usa veleni se necessario}
\statrow{Ombra}{1}{4}{Si muove negli ambienti sociali senza farsi notare}
\statrow{Combattimento}{0}{4}{Evita lo scontro diretto}
\statrow{Forza Bruta}{0}{4}{Non è la sua specialità}

\attrblock{doria}{RESOLVE}
\statrow{Percezione}{2}{4}{Legge le emozioni, sente le dinamiche nascoste}
\statrow{Comando}{2}{4}{Guida con la logica, non con il volume}
\statrow{Relazioni}{3}{4}{Eccelle nei salotti, nelle trattative, nelle alleanze}
\statrow{Persuasione}{3}{4}{Convince razionalmente o emotivamente a scelta}

\abilitybox{doria}{Due Mosse Avanti}{%
  Quando Bianca pianifica un'azione (non la esegue direttamente), può dichiarare
  una \textbf{Contingenza}: se l'azione fallisce, la contingenza si attiva
  automaticamente senza tiro.%
}{Va dichiarata prima del tiro dell'altro PG. Una volta per colpo.}

\subsection*{Legami}
\bondentry{Ottaviano}{Lo amo come un padre. Ma sta invecchiando — alcune decisioni le prendo io ormai.}
\bondentry{Ser Fulvio}{Gli voglio bene come a uno zio. Anche per questo mi spaventa: so che mente, ma non ancora su cosa.}
\bondentry{Fra' Matteo}{L'unico che non mi sottovaluta. Forse perché è abbastanza vecchio da non importargli di avere torto.}
\bondentry{Giacomo}{Utile. Pericoloso. Non mi fido di lui, ma ha accesso a canali che noi non abbiamo.}

\subsection*{Equipaggiamento}
\begin{itemize}
  \item Taccuino cifrato con la pianificazione completa del Colpo Riservato
  \item Lettera firmata da un senatore veneziano — un favore da riscuotere
  \item Abiti da uomo per muoversi nei quartieri chiusi alle donne
  \item Piccola fiala di sonnifero non letale — «per le emergenze»
\end{itemize}

\begin{secretbox}
% @rpg-schema: pg:BiancaDoria rpg:hasSecret
%   "Ha stabilito un contatto con i Grimaldi per una via di fuga della famiglia se la rivolta fallisce. Il prezzo non è ancora stato accettato." .
Bianca ha già stabilito un contatto con i Grimaldi per assicurare una via
di fuga alla famiglia se la rivolta fallisce. L'accordo prevede che i Doria
possano rifugiarsi a Monaco. Il prezzo non lo ha ancora accettato — ma i
Grimaldi aspettano la risposta.

\medskip
\textbf{\textcolor{SecretRed}{Agenda:}} Garantire la sopravvivenza del casato a lungo termine.
Se il Colpo Riservato mette in pericolo più persone di quante ne salvi,
proporrà un piano alternativo — anche contro il volere di Ottaviano.
\end{secretbox}

\newpage

% ---------------------------------------------------------------------------
% PG 4: Fra' Matteo da Recco
% @rpg-schema: pg:FraMatteo a fitd:Character ;
%   rdfs:label "Fra' Matteo da Recco" ;
%   fitd:crew genova1507:CasaDoria ;
%   rpg:role "Cappellano / La Coscienza" .
% ---------------------------------------------------------------------------

\charheader{doria}{Fra' Matteo da Recco}{Il Cappellano / La Coscienza}%
  {Dio perdona. Io tengo il conto.}

\begin{multicols}{2}
\subsection*{Background \& Motivazione}
Sessantun anni, domenicano, cappellano delle galee Doria da trent'anni.
Ha visto abbastanza battaglie navali da non farsi più illusioni
sulla natura umana. Crede ancora in Dio — non è sicuro di credere nell'umanità.

\textbf{\textcolor{DoriaBlue}{Motivazione:}} \textit{Un popolo cristiano non può essere
governato da potenze straniere senza il consenso della Chiesa.
Il Papa non ha benedetto la dominazione francese. La rivolta è giusta.}
\columnbreak
\subsection*{Stress \& Trauma}
\clocktrack[DoriaBlue]{STRESS}{0}{9}
\textit{\small Traumi: □ Freddo \ □ Duro \ □ Ossessionato \ □ Paranoico \ □ Spezzato}
\end{multicols}

\attrblock{doria}{INSIGHT}
\statrow{Caccia}{0}{4}{Non è il suo campo}
\statrow{Studio}{3}{4}{Teologia, medicina, storia ecclesiastica, latino e greco}
\statrow{Osserva}{2}{4}{Osserva le persone come se stesse leggendo le loro anime}
\statrow{Ingegno}{1}{4}{Preparazione di medicamenti, piccola chimica}

\attrblock{doria}{PROWESS}
\statrow{Finezza}{1}{4}{Mani ferme da chirurgo di bordo}
\statrow{Ombra}{0}{4}{Non è la sua specialità}
\statrow{Combattimento}{0}{4}{Non combatte. Non più.}
\statrow{Forza Bruta}{0}{4}{No.}

\attrblock{doria}{RESOLVE}
\statrow{Percezione}{3}{4}{Sente le dinamiche spirituali di un gruppo}
\statrow{Comando}{2}{4}{L'autorità morale può essere più potente di quella militare}
\statrow{Relazioni}{3}{4}{Confraternite, ordini religiosi, bassi ranghi del clero}
\statrow{Persuasione}{2}{4}{La voce da predicatore. Usata raramente, efficace.}

\abilitybox{doria}{Confessione \& Assoluzione}{%
  Quando un PG del tavolo sta per subire Trauma, Fra' Matteo può intervenire.
  Il Trauma diventa un semplice punto di Stress — ma il PG deve confessare
  a Fra' Matteo cosa lo ha portato a quel punto.%
}{Una volta per PG per sessione. Fra' Matteo deve essere fisicamente presente.}

\subsection*{Legami}
\bondentry{Ottaviano}{Lo conosco da quando aveva i capelli neri. Ha portato fardelli che non avrebbe dovuto.}
\bondentry{Ser Fulvio}{Non mi fido di lui. Mi confessa i peccati minori con tale regolarità che mi chiedo cosa stia nascondendo dietro di loro.}
\bondentry{Bianca}{Lei è il futuro di questa famiglia. Prego per lei ogni mattina.}
\bondentry{Giacomo}{Un peccatore allegro. Li preferisco ai santi tristi.}

\subsection*{Equipaggiamento}
\begin{itemize}
  \item Breviario con fogli sciolti nascosti dentro — messaggi cifrati
  \item Kit medico da campo: bende, alcol, ago e filo robusto
  \item Lettera del vescovo che lo autorizza a muoversi in tutti i sestieri
  \item Rosario di legno di ulivo — l'unica cosa rimasta del suo convento
\end{itemize}

\begin{secretbox}
% @rpg-schema: pg:FraMatteo rpg:hasSecret
%   "Conosce il debito segreto di Ottaviano sulle galee, appreso in confessione da un banchiere. Il sigillo confessionale lo paralizza." .
Fra' Matteo sa del debito segreto di Ottaviano sulle galee — non dalla confessione
di Ottaviano, ma da quella del banchiere, che si è confessato con lui.
Il sigillo confessionale lo paralizza. Non può dirlo.
Ma potrebbe agire indirettamente per risolvere il problema.

\medskip
\textbf{\textcolor{SecretRed}{Agenda:}} Proteggere Ottaviano da se stesso. Se il vecchio
ammiraglio prende una decisione che lo distruggerà, Fra' Matteo farà di tutto
per impedirlo, anche contro gli ordini espliciti.
\end{secretbox}

\newpage

% ---------------------------------------------------------------------------
% PG 5: Giacomo detto il Moro
% @rpg-schema: pg:GiacomoMoro a fitd:Character ;
%   rdfs:label "Giacomo detto il Moro" ;
%   fitd:crew genova1507:CasaDoria ;
%   rpg:role "Contrabbandiere / L'Uomo di Strada" .
% ---------------------------------------------------------------------------

\charheader{doria}{Giacomo detto il Moro}{Il Contrabbandiere / L'Uomo di Strada}%
  {Non sono dei vostri. Ma per stasera, forse sì.}

\begin{multicols}{2}
\subsection*{Background \& Motivazione}
Trentasei anni, origini oscure (si dice madre genovese, padre saraceno —
lui non smentisce perché gli torna utile). Contrabbandiere, sensale,
e occasionalmente informatore per chiunque paghi.

\textbf{\textcolor{DoriaBlue}{Motivazione:}} \textit{Denaro, prima di tutto. Ma anche un
senso oscuro che Genova libera sarebbe meglio per la gente come lui:
quella senza casato, senza nome, senza garanzie.}
\columnbreak
\subsection*{Stress \& Trauma}
\clocktrack[DoriaBlue]{STRESS}{0}{9}
\textit{\small Traumi: □ Freddo \ □ Duro \ □ Ossessionato \ □ Paranoico \ □ Spezzato}
\end{multicols}

\attrblock{doria}{INSIGHT}
\statrow{Caccia}{3}{4}{Trova persone, merci, informazioni in qualsiasi porto}
\statrow{Studio}{1}{4}{Lingue (arabo, catalano, greco storto)}
\statrow{Osserva}{3}{4}{Sopravvissuto per trent'anni grazie a questo}
\statrow{Ingegno}{2}{4}{Barche, reti, nodi, meccanismi semplici}

\attrblock{doria}{PROWESS}
\statrow{Finezza}{2}{4}{Borseggio, scalata, movimenti precisi in spazi stretti}
\statrow{Ombra}{3}{4}{Caruggi, tetti, canali — ogni via non ufficiale}
\statrow{Combattimento}{2}{4}{Coltello, bastone, combattimento sporco}
\statrow{Forza Bruta}{1}{4}{Sfonda quando necessario}

\attrblock{doria}{RESOLVE}
\statrow{Percezione}{1}{4}{Istinto animale per il pericolo}
\statrow{Comando}{0}{4}{Non comanda. Non è nel suo vocabolario.}
\statrow{Relazioni}{2}{4}{Taverne, mercati neri, basse compagnie}
\statrow{Persuasione}{2}{4}{Mercanteggia, convince, bluffa}

\abilitybox{doria}{I Caruggi Sono Miei}{%
  Nel sestiere di Prè, nel porto vecchio, o nei \idx{caruggi} del centro,
  Giacomo non può essere sorpreso. Ha sempre una via di fuga.
  Se inseguito, può far sparire sé stesso e un'altra persona senza tiro.%
}{Sempre attiva nei quartieri indicati. Altrove richiede Tiro su Ombra.}

\subsection*{Legami}
% @rpg-schema: pg:GiacomoMoro fitd:hasBond
%   [ fitd:bondTarget pg:FulvioDoria ; fitd:bondText "Siamo cugini veri. Lo ricordo quando fa comodo a me." ] .
\bondentry{Ottaviano}{Mi rispetta come si rispetta uno strumento utile. Va bene così.}
\bondentry{Ser Fulvio}{Siamo cugini. Lui lo dimentica quando fa comodo. Lo ricordo io, quando fa comodo a me.}
\bondentry{Bianca}{L'unica Doria che mi parla come a una persona. Non so se mi piaccia o mi spaventi di più.}
\bondentry{Fra' Matteo}{Mi confesso da lui perché è l'unico che non si sorprende di nulla.}

\subsection*{Equipaggiamento}
\begin{itemize}
  \item Barca a remi nascosta sotto il molo di Sampierdarena — solo lui sa dove
  \item Mappa mentale delle vie sotterranee sotto il porto (non scritta da nessuna parte)
  \item Coltello da caccia con manico in osso di balena
  \item Sacchetto con monete di cinque valute diverse — sempre pronto a partire
\end{itemize}

\begin{secretbox}
% @rpg-schema: pg:GiacomoMoro rpg:hasSecret
%   "Ha venduto agli Spinola informazioni sulla rotta del convoglio francese, centrale nel Colpo Riservato dei Doria." .
Giacomo ha già venduto informazioni sui Doria agli Spinola — non le più
pericolose, solo abbastanza per tenersi una porta aperta. Ma una delle
informazioni riguardava la rotta del convoglio francese — quella al centro
del Colpo Riservato. Se gli Spinola la usano, il Colpo è compromesso.

\medskip
\textbf{\textcolor{SecretRed}{Agenda:}} Sopravvivere e uscire da Genova con abbastanza denaro
da ricominciare altrove. L'unica cosa che potrebbe fermarlo è Fulvio —
se il cugino è nei guai, Giacomo resta.
\end{secretbox}

% =============================================================================
% I restanti 15 PG (Spinola, Fieschi, Grimaldi) seguono la stessa struttura.
% Per brevità del sorgente, sono inclusi nei file separati che seguono.
% Il file main.tex li include come \input{characters/spinola}, ecc.
% =============================================================================

% NOTE PER IL COMPILATORE:
% I personaggi degli Spinola, Fieschi e Grimaldi seguono
% esattamente la stessa struttura di questo capitolo.
% I dati completi sono nei file:
%   characters/spinola_chars.tex
%   characters/fieschi_chars.tex
%   characters/grimaldi_chars.tex
% inclusi dal main.tex.
